\ProvidesFile{hw_2}[Homework 2]

\begin{center}
Linear Algebra 2. Harbour.Space University

Graded homework 2. Part 1: mathematics.
\end{center}

Write up solutions to the problems below. Take photos or scan your solutions and combine them into {\bf a single PDF file}. Submit in Google Classroom by the deadline.

\begin{enumerate}[1. ]
\item Using linear operator theory (eigenvectors, eigenvalues, diagonal form), find the closed formula for the $n$-th Fibonacci number. Recall that the Fibonacci numbers are defined by
$$
F_0 = 0, \quad F_1 = 1, \quad F_n = F_{n-1} + F_{n-2} \quad \text{for } n \geq 2.
$$

\item Find the Jordan normal form and the transition matrix for the matrix
$$
\begin{pmatrix}
        2 & 0 & 1 & 1 & 0 & 0 \\
        1 & 1 & 1 & 0 & 1 & 0 \\
        -1 & 0 & 0 & 0 & 0 & 1 \\
        0 & 0 & 0 & 2 & 0 & 1 \\
        0 & 0 & 0 & 1 & 1 & 1 \\
        0 & 0 & 0 & -1 & 0 & 0
\end{pmatrix}.
$$
(You may use a calculator for arithmetic operations.)

Hints: use Problem 2.12 for determinant calculation. After you find the Jordan normal form $J$ and the transition matrix $C$, check (using any available software) that $J = C^{-1}AC$.

\item Note that if a linear operator $\mathcal{R}$ on $\R^2$ (two-dimensional real space --- a plane) has a matrix
$$
R' = \begin{pmatrix}
    1 & 0\\
    0 & -1\\
\end{pmatrix}
$$
in some 'good' basis, then this operator is a reflection and its eigenvector $v_1$ is a reflection axis.

\begin{enumerate}[a) ]
\item Find its second eigenvector in the basis  $v_1, v_2$.

\item Fix some basis $e_1, e_2$ and consider a reflection across the line through the vector $v = (x, y)$ in this basis. Then the transition matrix from this basis to the 'good' basis where the reflection has matrix $R'$ is 
$$
C = \begin{pmatrix}
    x & -y\\
    y & x\\
\end{pmatrix}.
$$
Prove this.

\item Find the matrix $R$ of the reflection across the line through the vector $v = (x, y)$ in the basis $e_1, e_2$.

\item Prove that a composition of two reflections is a rotation by an angle twice the angle between the reflection axes (counted counter-clockwise from the first axis to the second axis). In other words,
$$
R_2 R_1
= \begin{pmatrix}
    \cos 2\theta & -\sin 2\theta\\
    \sin 2\theta & \cos 2\theta\\
\end{pmatrix}
$$
where $R_1$ and $R_2$ are the matrices of the reflections across the lines through the vectors $(x_1, y_1)$ and $(x_2, y_2)$ respectively, and $\theta$ is the angle from vector $(x_1, y_1)$ to vector $(x_2, y_2)$. Note that composition of linear operators is applied from right to left; it acts on a vector as follows:
$$
R_2 R_1
\begin{pmatrix}
    \alpha\\
    \beta\\
\end{pmatrix}
= \begin{pmatrix}
    \cos 2\theta & -\sin 2\theta\\
    \sin 2\theta & \cos 2\theta\\
\end{pmatrix}
\begin{pmatrix}
    \alpha\\
    \beta\\
\end{pmatrix}.
$$

\item Consider a rotation by an angle $\varphi$, i. e. a linear operator with the matrix
$$
B =\begin{pmatrix}
    \cos \varphi & -\sin \varphi\\
    \sin \varphi & \cos \varphi\\
\end{pmatrix}
$$
as a real linear operator on $\R^2$ (two-dimensional real space) and as a complex linear operator on $\C^2$ (two-dimensional complex space). 

\begin{enumerate}[(1)]
\item Is it diagonalisable as a real operator? (Hint: if it was, it would have two real eigenvectors. What would it mean for the geometric properties of the operator?)
\item Is it diagonalisable as a complex operator? Find its Jordan normal form. Explain how its complex Jordan normal form relates to its geometric properties.
\end{enumerate}
\end{enumerate}

\item Recall that a Taylor series for an exponential function $e^x$ is
$$
e^x = \sum_{n=0}^{\infty} \frac{x^n}{n!}.
$$

Find the exponent of a Jordan block $J_n(\lambda)$, where exponent of a square matrix $A$ is defined by the formula
$$
e^A = \sum_{n=0}^{\infty} \frac{A^n}{n!}.
$$
(You may use the formula for $J_n(\lambda)^k$ from class work without proof.)


\end{enumerate}